\subsection{Ортогональные матрицы как способ задания ориентации тела}
Поворот твердого тела может быть задан с помощью ортогонального отображения трехмерного евклидова пространства в себя. Пусть есть два вектора: $\vec{r}=\begin{pmatrix}x\\y\\z\\ \end{pmatrix}$ и $\vec{r'}=\begin{pmatrix}x'\\y'\\z'\\ \end{pmatrix}$, в который отображается $\vec{r}$ с помощью преобразования поворота, которое задается матрицей U ($\vec{r'} = U\vec{r}$).\\
Из определения поворота $\vec{r}\cdot\vec{r} = \vec{r'}\cdot\vec{r'}$ $\Longrightarrow$
\[
\vec{r}\cdot\vec{r} = U\vec{r} \cdot U\vec{r} = \vec{r} \cdot U^TU\vec{r} \ \Longrightarrow \ U^TU = E
\]
Т.е. матрица $U$ ортогональна.\\
\underline{Некоторые факты про ортогональные матрицы:}\\
\begin{enumerate}
    \item $det~U = \pm 1$.
    \item $U^{-1} = U^T$.
    \item Произведение ортогональных матриц --- ортогональная матрица.
    \item Группа с произведением --- $O(3)$, подгруппа с $det = 1$ --- $SO(3)$.
    \item $SO(3)$ можно поставить во взаимно однозначное соответствие с поворотом.
    \item $SO(3)$ параметризует конфигурационное многообразие твердого тела с неподвижной точкой.
\end{enumerate}
Группа $SO(3)$ представляет собой конфигурационное многообразие твердого тела с одной неподвижной точкой.\\

Рассмотрим поворот тела на угол $\phi$ относительно осей:
\hypertarget{turns}{}
\[
x: \tab U = \begin{pmatrix}
    1 & 0 & 0\\
    0 & \cos{\phi} & -\sin{\phi}\\
    0 & \sin{\phi} & \cos{\phi}\\
\end{pmatrix}
\]
\[
y: \tab U = \begin{pmatrix}
    \cos{\phi} & 0 & \sin{\phi}\\
    0 & 1 & 0\\
    -\sin{\phi} & 0 & \cos{\phi}\\
\end{pmatrix}
\]
\[
z: \tab U = \begin{pmatrix}
    cos{\phi} & \sin{\phi} & 0\\
    sin{\phi} & \cos{\phi} & 0\\
    0 & 0 & 1\\
\end{pmatrix}
\]

Собственные векторы ортогональных преобразований удовлетворяют условию $U\vec{r} = \lambda\vec{r}$, где $\lambda$ --- число. Т.к. ортогональное преобразование не меняет длину вектора, $|\lambda| = 1$. Собственные числа находятся из характеристического уравнения, которое можно переписать в виде $\lambda^3 - \lambda^2tr~U + \lambda ~tr~U - 1 = (\lambda - 1)(\lambda^2 + \lambda - \lambda~tr~U + 1) = 0$, где след матрицы $tr~U = u_{11} + u_{22} + u_{33}$. Уравнение имеет корни $\lambda_1 = 1$, 
\[
\lambda_{2,~3} = \dfrac{tr~U - 1}{2} \pm \sqrt{\dfrac{(tr~U - 1)^2}{4} - 1}
\]
А т.к. модуль определителя матрицы равен 1, $(tr~U)^2 - 2~|tr~U| - 3 < 0 $ $\Longrightarrow$ $\lambda_{2,~3} = \cos{\phi} \pm i\sin{\phi}$ $\Longrightarrow$ $\cos{\phi} = \dfrac{tr~U - 1}{2}$.\\
Найдем собственные векторы $\vec{r}_1$, $\vec{r}_2$ и $\vec{r}_3$. Для этого решим три однородные системы $U\vec{r}_k = \lambda_k\vec{r}_k$. Так, тривиально находится первый вектор: $\vec{r}_1 = \vec{r}$. А для $\lambda_{2,~3}$ векторы имеют вид $\vec{r}_{2,~3} = \vec{p} \mp i\vec{q}$, где $\vec{p}$ и $\vec{q}$ --- вещественные векторы.
\begin{lemma}
    Тройка ($\vec{r}, \vec{p}, \vec{q}$) --- правый ортогональный базис.
\end{lemma}
\begin{proof}
    С одной стороны
    \[
    (U\vec{r_1})^T\cdot \vec{r_2} = \vec{r_1}^T\cdot \vec{r_2}
    \]
    А с другой
    \[
    (U\vec{r_1})^T\cdot \vec{r_2} = \vec{r_1}^T \cdot U^T \cdot \vec{r_2} = \lambda_3 \vec{r_1}^T \cdot \vec{r_2},
    \]
    т.к. $\vec{r_2}$ --- собственный вектор матрицы $U^T$. Сравнивая получим, что если $\lambda_3 \neq 1$, то $\vec{r_1}^T \cdot \vec{r_2} = 0$ $\Longrightarrow$ $\vec{r}^T\cdot\vec{p} = \vec{r}^T\cdot\vec{q} = 0$.\\
    Аналогично предыдущему шагу, для $(U\vec{r_2})^T\cdot \vec{r_2}$ получаем, что если $\lambda_2 \neq \lambda_3$, то $\vec{r_2}^T\cdot\vec{r_2} = 0$ $\Longrightarrow$ $\vec{p}^2 = \vec{q}^2$ и $\vec{p}^T\cdot \vec{q} = 0$.\\

    Случай $\lambda_1 = \lambda_2 = \lambda_3 = 1$ --- тождественное преобразование $U = E$, для которого все векторы собственные $\Longrightarrow$ найдутся взаимно ортогональные.\\

    Для того, чтобы показать, что это правая тройка, найдем векторы $\vec{r_2},~\vec{r_3}$:
    \[
    \vec{r_2} = \begin{pmatrix}
        0\\1\\-i\\
    \end{pmatrix}, \tab \vec{r_3} = \begin{pmatrix}
        0\\1\\i\\
    \end{pmatrix}, \tab \vec{p} = \begin{pmatrix}
        0\\1\\0\\
    \end{pmatrix}, \tab \vec{q} = \begin{pmatrix}
        0\\0\\1\\
    \end{pmatrix}
    \]
    Понятно, что тройка ($\vec{r}, \vec{p}, \vec{q}$) --- правая.
\end{proof}

Следующая теорема является следствием выкладок выше.
\begin{theorem}[Теорема Эйлера о поворотах]
    Любое перемещение твердого тела с одной неподвижной точкой может быть заменено плоским поворотом вокруг некоторой оси $\vec{r_1}$ на некоторый угол ($\cos{\phi} = \dfrac{tr~U - 1}{2}$).
\end{theorem}

Еще одно свойство матриц $U\in SO(3)$: любое ортогональное преобразование пространства эквивалентно повороту пространства вокруг $\vec{r}$ на угол $\phi$. Для доказательства достаточно рассмотреть уравнения $U\vec{r_k} = \lambda_k\vec{r_k}$ в вещественной форме:
\[
\begin{cases}
    U\vec{r} = \vec{r}\\
    U\vec{p} = \vec{p}\cos{\phi} + \vec{q}\sin{\phi}\\
    U\vec{p} = -\vec{p}\sin{\phi} + \vec{q}\cos{\phi}\\
\end{cases}
\]
Получим
\[
U = \begin{pmatrix}
    1 & 0 & 0\\
    0 & \cos{\phi} & -\sin{\phi}\\
    0 & \sin{\phi} & \cos{\phi}\\
\end{pmatrix}
\]
